% !TeX root = ../main.tex
\section{Short-strip detector}
\label{sec:short-strips}
The intermediate tracker uses rectangular two-coordinate sensors, referred to
here as strixels. The selected layout combines four tilted barrel layers with
six petal discs at each end. It retains a common module contract for barrel and
endcap, with the fine coordinate transverse to the track's radial direction and
the coarse coordinate longitudinal in the barrel or radial in the endcap.

\subsection{Sensor, electronics and thermal stack}
Each module has a \qty{48}{\milli\metre} by \qty{96}{\milli\metre} active area
in a \qty{0.2}{\milli\metre} silicon sensor, with a \qty{0.5}{\milli\metre}
guard on every edge. Its \qty{75}{\micro\metre} by \qty{0.5}{\milli\metre}
grid contains $640\times192=\num{122880}$ nominal cells. Eight
\qty{23.4}{\milli\metre}-square, \qty{0.15}{\milli\metre}-thick readout tiles
occupy a $2\times4$ array at \qty{24}{\milli\metre} centre pitch. This is a
footprint hypothesis: redistribution, tile gaps, channel compatibility and
readout power require a real electronics design. Cell counts do not establish
live-channel coverage across the tile boundaries.

The modeled module includes a \qty{25}{\micro\metre} bump envelope, represented
by copper and air fractions, followed by \qty{0.10}{\milli\metre} adhesive,
\qty{0.20}{\milli\metre} graphite, \qty{25}{\micro\metre} insulation,
\qty{0.10}{\milli\metre} adhesive and \qty{0.20}{\milli\metre} CFRP backing.
Two \qty{3}{\milli\metre} axial flex ledges enlarge the occupied outline to
\qty{49}{\milli\metre} by \qty{103}{\milli\metre}. Their bridges and tails are
included in clearance checks. The graphite and local pickups couple the module
to the cooled sandwich. Neither the assumed contact resistance nor the carbon
conductivity is a qualified cold interface.

\subsection{Tilted barrel placement}
The stave anchor radii are $260,340,480,660\,\unit{\milli\metre}$, with
$44,56,80,108$ staves respectively. Each carries 28 modules, giving 288 staves
and \num{8064} sensors. The modules have symmetric endpoint centres at
$z=\pm\qty{1152}{\milli\metre}$ and \qty{85.333}{\milli\metre} row pitch;
their axial active overlap is \qty{10.667}{\milli\metre}. Alternating rows
are lifted by \qty{1.5}{\milli\metre} to separate complete occupied modules.
The endpoint active edges reach the nominal $|z|=\qty{1200}{\milli\metre}$.
The tilt angle has been optimized to $+\qty{15}{\degree}$ about global $z$.

\begin{figure}[tbp]
\centering\includegraphics[width=\linewidth]{short-strip-barrel}
\caption{Selected short-strip barrel. Left: both row elevations are projected
for every stave; labels give layer populations. Right: the innermost cold-stack
projection shows equal-radius shingling, embedded tubes and the bearing contact
on a bare skin land. It superposes different longitudinal rows, rather than
showing one physical $z$ cut. The common azimuthal tilt is \qty{15}{\degree}.}
\label{fig:short-barrel}
\end{figure}

\subsection{Structural stave, bearings and barrel routes}
The \qty{52}{\milli\metre}-wide cold structure extends to
$|z|=\qty{1210}{\milli\metre}$. It contains a \qty{5}{\milli\metre} conductive
carbon-foam core, \qty{0.10}{\milli\metre} core adhesive and
\qty{0.15}{\milli\metre} CFRP skins on each face. The skins carry bending load;
the core separates them, transfers shear and houses the cooling. Module pickups
follow each row's elevation. Both pipes and their coolant volumes are explicit,
and their full outer envelopes are subtracted from the core.

Nine CFRP bearing rings per layer are centred at
$z=0,\pm250,\pm550,\pm850,\pm1150\,\unit{\milli\metre}$. Each is
\qty{8}{\milli\metre} wide axially and \qty{3}{\milli\metre} radially.
A beveled \qty{2}{\milli\metre}-wide web at local
$u=\qty{14.5}{\milli\metre}$ joins the bare back-skin land to the ring tangent
plane, leaving \qty{0.5}{\milli\metre} to the adjacent insulation and signal
bus.

\begin{figure}[tbp]
\centering\includegraphics[width=\linewidth]{short-strip-stave}
\caption{Short-strip stave structure and bearing schedule. Top: local layer
positions and embedded tube sections; low and raised sensors are superposed.
Thin module films and complete pickup details are omitted from this explanatory
view. Bottom: the nine selected bearing stations and the two independently
served half-staves. End boards lie beyond the active rows.}
\label{fig:short-stave}
\end{figure}

Each signed half-stave has one titanium U-loop with two straight legs at
$u=\pm\qty{12}{\milli\metre}$, tube outside diameter
\qty{2.5}{\milli\metre}, wall \qty{0.14}{\milli\metre}, and centre
$w=-\qty{4.05}{\milli\metre}$. Straight runs extend from
$|z|=16$ to \qty{1210}{\milli\metre}; the central return has
\qty{12}{\milli\metre} centre radius. Positive and negative circuits are
hydraulically separate. The backside LV carrier holds two
\qty{12}{\milli\metre} by \qty{0.20}{\milli\metre} copper rails on
\qty{0.15}{\milli\metre} polyimide. HV, control and data use a separate patterned
flex. Module tails route around the cold structure to these insulated buses.

End-board allocations are \qty{52}{\milli\metre} by \qty{30}{\milli\metre} by
\qty{8}{\milli\metre}, at $z=\pm\qty{1230}{\milli\metre}$. Each half-stave's
14 modules require two harnesses under the 12-module limit. Each end therefore
presents 576 harnesses and 288 cooling loops, grouped into twelve sector
collectors at $r=244$--$750\,\unit{\milli\metre}$,
$|z|=1245$--$1310\,\unit{\milli\metre}$, before the fixed outer trunk.
The gross harness fixture uses \qty{13.4}{\milli\metre} power and
\qty{3.6}{\milli\metre} fibre envelopes. These dimensions are capacity
comparators, not a manufactured cable specification.

\subsection{Endcap discs and removable petals}
Six discs per end use the datums in \cref{tab:strip-discs}. The nominal active
annulus spans $r=243.5$--$670.06\,\unit{\milli\metre}$. Twelve
\qty{30}{\degree} petals divide each disc into removable cold assemblies.
Five radial rings carry $48,60,72,84,96$ modules, or 30 per petal and 360 per disc;
the two endcaps contain \num{4320} sensors. Fine $U$ is tangential, coarse $V$
points inward radially, and $N$ points outward from the interaction region.
Alternating ring and azimuth elevations give four offsets
$0,1.5,3,4.5\,\unit{\milli\metre}$. The rectangular bodies overlap in
projection while their full passive stacks remain separated in depth.

\begin{table}[tbp]
\centering\small
\caption{Positive support datums for the strip discs. The opposite end uses
the corresponding signed placements. Values are rounded for display. The first
short- and long-strip datums include the barrel-collector clearance amendments.}
\label{tab:strip-discs}
\begin{tabular}{@{}rrr@{}}\toprule
Disc & Short strips $|z|$ [mm] & Long strips $|z|$ [mm]\\\midrule
\csname @@input\endcsname generated/strip-disc-rows.tex
\bottomrule\end{tabular}
\end{table}

\begin{figure}[tbp]
\centering\includegraphics[width=\linewidth]{short-strip-endcap}
\caption{Short-strip endcap, with all four elevations projected. Left: the
first positive disc's five rectangular module rings. Right: one petal and its
embedded serpentine cooling centreline. The two backside bus routes are shown
schematically at their selected azimuths. The outer board, fanout transitions
and individual connectors require detailed routing review.}
\label{fig:short-endcap}
\end{figure}

The petal spans $r=236$--$685\,\unit{\milli\metre}$ with a
\qty{0.6}{\milli\metre} inner seam gap. Its \qty{5}{\milli\metre} foam,
\qty{0.10}{\milli\metre} glue and \qty{0.15}{\milli\metre} skins form the
cooled structural plate. Two \qty{6}{\milli\metre} by \qty{60}{\milli\metre}
graphite pickups per module follow the elevations; their small thermal area is
included in the adverse screen. Inner and outer CFRP frames occupy
$r=236$--$250$ and $674$--$700\,\unit{\milli\metre}$.

One serpentine loop per petal runs on ten circular tracks at the five module
radii plus or minus \qty{20}{\milli\metre}, with radial U-turns between tracks.
The feed and return approach the rim at $29$ and \qty{28}{\degree}, through
\qty{10}{\milli\metre} entry fillets. The tube dimensions match the barrel.
LV and HV/data buses remain separate on the back; front fanout cells preserve
flex constituents and approximate length above the modules, without representing
every connector, via or vertical transition. The
\qty{52}{\milli\metre} by \qty{20}{\milli\metre} by \qty{8}{\milli\metre}
outer board sits at $r=\qty{695}{\milli\metre}$. Each petal supplies three
harnesses and one loop to its $r=685$--$710\,\unit{\milli\metre}$ outer fan,
then to the shared trunk. The endcaps add 216 harnesses and 72 loops per end,
giving combined barrel/endcap totals of 792 harnesses and 360 loops per end.
